\chapter{Cognition and Memory}
\index{memory}\index{cognition}

\section{What Kind of Object Is a Memory?}

From the framework's perspective, memory is a compression operator on experiential trajectories. The neural mechanisms, storage architecture, and retrieval dynamics are implementation details. What matters is the compression: what equivalence classes it induces on experiential trajectories, what queries it supports, and whether it is aligned with the stability structure of experiential dependence.

Before asking whether memory is stable or sufficient, the framework must ask: what is the trajectory space?

\section{Candidate Trajectory Spaces}

Three candidates arise:

\textbf{Episodic.} A trajectory is a temporally ordered sequence of experiences. Counterfactual question: what would I believe had experience $X$ been replaced by $Y$?

\textbf{Inferential.} A trajectory is a sequence of belief states and inferential transitions. Counterfactual question: what would I conclude had I inferred differently at some earlier stage?

\textbf{Predictive.} A trajectory is a sequence of action-perception loops with continuous prediction-error updates. Connects to the free-energy principle \cite{Friston2010}.

\textbf{Diagnostic Question 1: Partially open.} Multiple candidate interaction orders exist (temporal, associative, causal depth); none is established as canonical.

\section{Stability: Diagnostic Question 2}

The forgetting curve \cite{Ebbinghaus1885} suggests $\sigma_r \sim r^{-k}$ under temporal interaction order---polynomial decay, consistent with a summable stability sequence. Interference effects complicate the picture; traumatic experiences may exhibit stability failure for specific query families.

\textbf{Diagnostic Question 2: Conditional.} Evidence suggests polynomial decay; exponential stability unestablished; trauma may produce stability failure.

\section{Compression and Failure Modes}

Memory induces an equivalence relation on trajectories. Candidate relations: semantic equivalence, narrative equivalence, emotional equivalence, predictive equivalence.

The failure-mode taxonomy maps cleanly onto recognized phenomena:
\begin{itemize}
  \item \textbf{False memories} $\leftrightarrow$ Type I degeneracy
  \item \textbf{Amnesia} $\leftrightarrow$ sufficiency failure
  \item \textbf{Overgeneralization} $\leftrightarrow$ structural misalignment
  \item \textbf{Intrusive memories} $\leftrightarrow$ stability failure (trauma)
\end{itemize}

\section{Diagnostic Audit and Open Obligations}

\begin{center}\small
\begin{tabular}{lll}
\toprule
Question & Status & Notes \\
\midrule
Decomposition exact? & Partially open & No canonical interaction order \\
Stability summable? & Conditional & Polynomial decay suggested \\
Compression sufficient? & Conditionally yes & Equivalence relation unknown \\
Depths aligned? & Unknown & Depends on Q1, Q2 \\
Degeneracy avoided? & Active concern & False memories suggest Type I \\
\bottomrule
\end{tabular}
\end{center}

Open obligations: (1) determine canonical interaction order; (2) measure stability sequence; (3) characterize memory equivalence classes; (4) measure reconstruction conditioning; (5) identify stability structure of traumatic memory.
